\documentclass[11pt,a4paper]{article}

\usepackage[margin=1in]{geometry}
\usepackage[T2A]{fontenc}
\usepackage[utf8]{inputenc}
\usepackage[ukrainian]{babel}
\usepackage{booktabs}
\usepackage{longtable}
\usepackage{array}
\usepackage{hyperref}
\usepackage{setspace}
\usepackage{csquotes}
\usepackage[
  backend=biber,
  style=apa,
  sorting=nyt
]{biblatex}
\addbibresource{adhd_psychometric_note.bib}
\DeclareLanguageMapping{ukrainian}{english-apa}

\setstretch{1.08}
\hypersetup{
  unicode=true,
  colorlinks=true,
  linkcolor=black,
  urlcolor=blue,
  pdftitle={Що цей тест обґрунтовано може показувати},
  pdfauthor={Grendel}
}

\title{Що цей тест обґрунтовано може показувати\\
\large Коротка психометрична нотатка}
\author{Grendel evidensbasert psykologi AS}
\date{\today}

\begin{document}
\maketitle

\begin{abstract}
Цей текст підсумовує, що обґрунтовано можна стверджувати про те, що насправді вимірює коротка, позитивно сформульована шкала самозвіту щодо початку справ, зосередженості, організації та саморегуляції. Головний висновок полягає в тому, що тест не можна використовувати як діагностичний інструмент для ADHD, але він, схоже, вловлює широкий і узгоджений вимір регуляторного тертя, релевантний для труднощів, подібних до ADHD. Водночас моделі показують, що цей широкий вимір має певну внутрішню структуру. Найбільш виправдана інтерпретація, отже, полягає в тому, що тест вимірює спільний регуляторний фактор із двома частково розрізнюваними проявами: більш внутрішнім патерном, пов'язаним із залученням у завдання та утриманням думок на курсі, і більш зовнішнім патерном, пов'язаним із практичною надійністю та організацією повсякдення.
\end{abstract}

\section{Вступ}
Про короткі опитувальники легко говорити так, ніби вони або «вловлюють ADHD», або «не вловлюють ADHD». Більш зріле розуміння полягає в тому, що такі інструменти зазвичай вимірюють патерн самозвітованих труднощів, який може бути більш або менш типовим для ADHD, не даючи при цьому підстав для діагнозу. Цей тест складається з десяти тверджень, сформульованих як твердження про функціонування, а не як твердження про симптоми. Тож він не запитує прямо, чи має людина ADHD, а про те, наскільки легко чи важко братися до справ, утримувати зосередженість, пам'ятати про домовленості, організовувати час, тримати думки на курсі та функціонувати без надзвичайно великої зовнішньої структури.

Мета цієї нотатки --- описати, що результати насправді підтверджують, і, що не менш важливо, чого вони не підтверджують. Амбіція полягає не в тому, щоб зробити тест більш певним, ніж він є, а в тому, щоб сформулювати точну проміжну позицію: тест, схоже, вимірює щось реальне та психометрично узгоджене, але це «щось» найкраще розуміти як патерн регуляторного тертя, релевантний для ADHD, а не як діагностичну відповідь.

\section{Дані та скринінг}
Поточний аналіз ґрунтувався на $N = 309$ необроблених відповідях. Набір даних містив практично лише відповіді букмолом, з дуже невеликою кількістю відповідей нюношком, англійською та квенською. У цьому запуску швидкі дублікати не видалялися, натомість було виключено 28 так званих рівних серединних відповідей, тобто відповідей, у яких на всі десять тверджень було дано категорію 3. Отже, аналізи оцінювалися на 281 відповіді.

Новіший запуск ренормування, використаний для оновлення продакшн-моделі, дав схожий, але дещо більший матеріал. Там було отримано 390 необроблених відповідей, з яких 346 збережено після фільтрації за якістю. Цей новіший запуск використовується тут лише як підтримка загальної інтерпретації, а не як заміна основного аналізу.

\section{Програмне забезпечення та схема аналізу}
Аналізи проводилися в \texttt{R} \parencite{RCore2026}. Центральним пакетом був \texttt{mirt} \parencite{Chalmers2012mirt}, який використовувався кількома способами:

\begin{itemize}
\item \texttt{mirt(..., itemtype = "graded")} використовувався для оцінювання graded response-моделей з одним фактором, двома факторами та біфакторною структурою.
\item \texttt{fscores()} використовувався для обчислення індивідуальних оцінок методами EAP і MAP, а також очікуваних латентних оцінок, обумовлених сумарним балом (\texttt{EAPsum}).
\item \texttt{anova()} використовувався для порівняння одно- та двофакторної моделей.
\item \texttt{M2()} використовувався як глобальний індекс відповідності для однофакторної моделі.
\item \texttt{itemfit()} використовувався для діагностики пунктів.
\item \texttt{residuals(type = "Q3")} використовувався для дослідження локальної залежності між пунктами.
\end{itemize}

Скрипт також завантажував \texttt{careless} \parencite{YentesWilhelm2023careless}, але в цьому запуску пакет не використовувався явно в самому оцінюванні. Скринінг якості в цьому аналізі складався на практиці з простої логіки, написаної в base-\texttt{R}: сортування за часом, видалення можливих швидких дублікатних патернів і виключення рівних серединних відповідей. Пакети \texttt{stats4} і \texttt{lattice} \parencite{Sarkar2008lattice} завантажувалися автоматично як залежності \texttt{mirt}; для \texttt{stats4} використовується стандартне цитування \texttt{R} \parencite{RCore2026}.

\section{Основні результати}
\subsection{Однофакторна модель}
Однофакторна модель дала чіткий і однорідний патерн. Факторні навантаження лежали між 0.659 і 0.825, і всі пункти мали помірні або високі спільності. Сума навантажень становила 5.682, а частка поясненої дисперсії --- 0.568. Надійність була високою як у вигляді $\alpha = 0.904$, так і як факторна надійність $r_{xx,F1} = 0.901$. Стандартна похибка вимірювання становила 3.066.

Глобальна відповідність однофакторної моделі в цьому запуску була доброю. Тест M2 дав $\chi^2(40)=16.177$, $p=1.00$, а структура залишків була спокійною. Залишки Q3 мали максимум 0.169 із середнім близько $-0.095$, що свідчить проти серйозної локальної залежності між пунктами. На цьому рівні тест, отже, поводиться як добре функціональна коротка шкала з одним чітким головним виміром.

\subsection{Двофакторна модель}
Коли було оцінено двофакторну модель, вона покращила відповідність порівняно з однофакторною. Порівняння дало $\Delta \chi^2 = 68.12$ з 9 ступенями свободи та $p < .001$. Це означає, що матеріал не є ідеально одновимірним у строгому статистичному сенсі. Питання, однак, у тому, що ця додаткова структура представляє.

У попередніх запусках могло здаватися, що ця додаткова структура частково стосувалася дефіциту уваги, а частково --- більш вільної індивідуальної варіації чи методичних ефектів. Новіше ротоване двофакторне рішення, оцінене на оновленому та відфільтрованому матеріалі, видається дещо більш інтерпретовним. Там пункти 1, 4, 5, 9 і 10 найсильніше навантажувалися на один фактор, тоді як пункти 2, 3, 7 і 8 --- на інший. Пункт 6 показував більш змішану належність. Водночас два фактори були високо корельованими ($r = 0.75$), що дуже важливо: додаткова структура вказує не на дві незалежні риси, а на два близько споріднені прояви ширшого регуляторного фактора.

\subsection{Біфакторна модель}
Біфакторна модель дає найбільш інформативне цілісне розуміння. У цьому аналізі всі десять пунктів чітко навантажувалися на генеральний фактор $G$ з навантаженнями від 0.574 до 0.792. Водночас виявилися два слабші групові фактори. Перший груповий фактор був переважно пов'язаний із пунктами 1--5, тоді як другий найчіткіше пов'язаний із пунктом 6 і меншою мірою з пунктами 7--10.

Вирішальне спостереження, втім, полягає в тому, що генеральний фактор ніс більшу частину спільної дисперсії. Частка поясненої дисперсії (Proportion Var) становила 0.529 для $G$ проти 0.039 і 0.065 для двох групових факторів. Іншими словами: тест має більше структури, ніж повністю вловлює чиста однофакторна модель, але головне враження залишається таким, що домінує один широкий фактор.

\section{Інтерпретація}
Найбільш виправданим твердженням є, отже, те, що тест, схоже, вловлює широкий і наскрізний вимір регуляторного тертя, релевантний для труднощів, подібних до ADHD. Сильний генеральний фактор свідчить на користь того, що в матеріалі існує принаймні один спільний фактор, який проходить крізь пункти і збирає те, що багато клініцистів і користувачів інтуїтивно впізнають як труднощі з підтриманням повсякдення в стабільний і самокерований спосіб.

Спокусливо назвати це «фактором, який повторюється в усіх людей з ADHD». Емпірично дані так далеко не сягають. У цьому матеріалі немає ні діагностичних даних золотого стандарту, ні зовнішньої валідизації проти усталених інструментів ADHD. Те, що дані насправді підтверджують, --- це стриманіше формулювання: шкала, схоже, вловлює широкий і, ймовірно, центральний регуляторний фактор, релевантний для ADHD, щодо якого правдоподібно думати, що він повторюватиметься в різних презентаціях ADHD. Це важливий результат, але це не те саме, що довести універсальність фактора для всіх людей з ADHD.

Вторинна структура все ж цікава й змістовна. Вона більше не виглядає як чистий шум. У новішому рішенні один бік тесту постає більш пов'язаним із внутрішнім залученням у завдання: починати, тримати увагу зібраною, не залежати від кризової енергії, утримувати думки на місці та обходитися без великої зовнішньої підтримки. Інший бік постає більш пов'язаним із зовнішньою надійністю та практичною організацією: не губити речі, дотримуватися домовленостей, пам'ятати повідомлення та змушувати повсякденну логістику працювати. Це можна розуміти як два прояви однієї надрядної сфери труднощів: більш внутрішній і більш зовнішній бік регуляції.

\section{Що тест обґрунтовано може показувати}
На практичному рівні, отже, обґрунтовано сказати, що тест показує, наскільки самозвітоване функціонування людини схоже чи не схоже на релевантний для ADHD патерн регуляторних труднощів. Вищі бали вказують на меншу схожість із такими труднощами. Нижчі бали вказують на більшу схожість. Це можна використовувати педагогічно і як структурований відправний пункт для рефлексії, особливо тому, що тест видається коротким, послідовним і психометрично порівняно чистим.

Натомість необґрунтовано казати, що тест вирішує, чи має людина ADHD. Низький бал може бути сумісним з ADHD, але також і з низкою інших станів, як-от стрес, недосипання, депресія, тривожність, перевантаження чи інші форми когнітивного та емоційного виснаження. Відповідно, високий бал може бути сумісним із відсутністю суттєвих регуляторних труднощів, але також і з легким ADHD, компенсованими труднощами чи рівнем функціонування, який добре підтримується структурою, інтелектом, рутинами чи адаптацією середовища.

\section{Обмеження}
Є кілька очевидних обмежень. По-перше, аналіз ґрунтується на самозвіті. По-друге, матеріалу бракує зовнішньої валідизації проти клінічної діагностики чи усталених шкал ADHD. По-третє, мовний розподіл є перекошеним, тож поки немає підстав для серйозних мовно-специфічних норм чи аналізів інваріантності вимірювання. По-четверте, біфакторне та двофакторне рішення є діагностичними моделями в методичному сенсі; вони корисні для розуміння структури даних, але самі по собі не легітимізують тлумачення тесту як двох окремих клінічних підшкал.

\section{Висновок}
Найзріліший і науково найвиправданіший опис тесту полягає в тому, що він працює як короткий, нормований показник регуляторного тертя з чіткою релевантністю для труднощів, подібних до ADHD. Результати підтверджують існування сильного й наскрізного генерального фактора, тобто спільного виміру, який пов'язує початок справ, зосередженість, організацію, внутрішній спокій і стабільне доведення до кінця. Водночас і двофакторна, і біфакторна моделі вказують на те, що цей генеральний фактор має певну внутрішню структуру, в якій більш внутрішній, пов'язаний із завданнями та увагою бік можна відрізнити від більш зовнішнього, практичного й пов'язаного з надійністю боку.

Те, що тест, отже, обґрунтовано може показувати, --- це не «ADHD» як діагноз, а ступінь схожості між відповідями людини та узгодженим патерном релевантного для ADHD регуляторного тертя. Це важливий і корисний результат. Але це все ще результат на рівні функціонування, а не остаточна відповідь щодо причини, категорії чи клінічної ідентичності.

\printbibliography

\end{document}
